\documentclass[11pt,a4paper]{article}

\usepackage[utf8]{inputenc}
\usepackage[T1]{fontenc}
\usepackage{lmodern}
\usepackage{microtype}
\usepackage{amsmath,amssymb,amsthm}
\usepackage{physics}      % \ket, \bra, \braket
\usepackage{booktabs}
\usepackage{hyperref}
\usepackage{cleveref}
\usepackage{xcolor}
\usepackage{listings}
\usepackage{geometry}
\usepackage{authblk}
\usepackage{abstract}
\usepackage{titlesec}
\usepackage{parskip}
\usepackage{graphicx}
\usepackage{caption}

\geometry{margin=2.5cm}

\hypersetup{
  colorlinks=true,
  linkcolor=blue!60!black,
  citecolor=blue!60!black,
  urlcolor=blue!60!black,
}

\definecolor{codegreen}{rgb}{0.13,0.55,0.13}
\definecolor{codegray}{rgb}{0.45,0.45,0.45}
\definecolor{codeblue}{rgb}{0.0,0.3,0.7}
\definecolor{backcolor}{rgb}{0.96,0.96,0.98}

\lstdefinestyle{rust}{
  language=C,
  backgroundcolor=\color{backcolor},
  basicstyle=\footnotesize\ttfamily,
  keywordstyle=\color{codeblue}\bfseries,
  commentstyle=\color{codegray}\itshape,
  stringstyle=\color{codegreen},
  breaklines=true,
  frame=single,
  framesep=4pt,
  rulecolor=\color{gray!40},
  xleftmargin=1em,
  xrightmargin=1em,
}
\lstdefinestyle{python}{
  language=Python,
  backgroundcolor=\color{backcolor},
  basicstyle=\footnotesize\ttfamily,
  keywordstyle=\color{codeblue}\bfseries,
  commentstyle=\color{codegray}\itshape,
  stringstyle=\color{codegreen},
  breaklines=true,
  frame=single,
  framesep=4pt,
  rulecolor=\color{gray!40},
  xleftmargin=1em,
  xrightmargin=1em,
}

% ─────────────────────────────────────────────────────────────────────────────
\title{\textbf{Neutral Cross-Framework Benchmarking Reveals\\
       $R_z$ Sign Convention Divergence in Quantum Simulation Backends}}

\author{Cleiton Augusto Correa Bezerra}
\affil{CleitonForge Project \\ \texttt{augusto.cleiton@gmail.com}}
\date{July 2026}

% ─────────────────────────────────────────────────────────────────────────────
\begin{document}

\maketitle

\begin{abstract}
We present \textsc{CleitonForge}, a neutral benchmarking layer for quantum
circuit simulation implemented in Rust. The framework routes identical circuits
through a canonical intermediate representation (IR) to multiple independent
simulation backends, enabling direct cross-backend comparison without
framework-specific assumptions.

Benchmarking four backends---a native statevector implementation, \texttt{quantrs2},
\texttt{roqoqo}, and \texttt{q1tsim}---we identify a sign convention
disagreement in the $R_z$ rotation gate. The divergence produces
cross-backend fidelity $F = 0.0$ on QAOA and VQE circuits but is
\emph{invisible} to Clifford benchmarks and the IBM Quantum Volume test,
because Haar-random circuits average over all rotation angles.

We characterize the affected gate family ($R_z$, Phase, $CR_z$, $CP$, $U$,
$CU$), implement a convention normalization transpiler that restores fidelity
to unity, and discuss implications for variational quantum algorithm
portability. \textsc{CleitonForge} is open source (Apache 2.0) and available
at \url{https://github.com/cleitonaugusto/cleitonforge}.
\end{abstract}

\textbf{Keywords:} quantum simulation, benchmarking, convention normalization,
QAOA, statevector simulation, Rust, open source.

% ─────────────────────────────────────────────────────────────────────────────
\section{Introduction}

The quantum software ecosystem has produced several high-quality statevector
simulation libraries in recent years. Users writing variational quantum algorithms
(QAOA, VQE, ADAPT-VQE) often prototype on one framework and execute on another
or on real hardware. Implicit in this workflow is an assumption: that the same
circuit, expressed in the same notation, produces the same quantum state across
all simulation backends.

This assumption fails silently in the presence of gate convention
disagreements. Unlike bugs that produce incorrect results on \emph{every} input,
convention disagreements produce correct results on most circuits (Clifford
gates, random circuits) while silently diverging on exactly the circuits that
matter most in practice: parameterized circuits with specific, non-trivial
angles.

\textsc{CleitonForge} was designed to surface this class of discrepancy. By
maintaining strict neutrality---no built-in preference for any framework, all
backends fed the same canonical IR---we can attribute observed fidelity
differences to backend-level convention choices rather than circuit-encoding
artifacts.

% ─────────────────────────────────────────────────────────────────────────────
\section{CleitonForge Architecture}

\textsc{CleitonForge} is a Rust workspace with four crates: \texttt{cforge-core}
(canonical IR), \texttt{cforge-parser} (OpenQASM~2/3 parsers and the normalization
transpiler), \texttt{cforge-backends} (adapters for each simulator), and
\texttt{cforge-metrics} (cross-backend comparison). Python bindings are provided
via \texttt{cforge-py} (PyO3/maturin).

\subsection{Canonical intermediate representation}

Every backend adapter receives a \texttt{Circuit} value---a flat sequence of
typed \texttt{Operation} records, each carrying a \texttt{GateKind} enum, qubit
indices, and floating-point parameters. No framework-specific types cross the
adapter boundary. The input to every backend is bit-for-bit identical.

\begin{lstlisting}[style=rust, caption={Canonical IR excerpt (\texttt{cforge-core}).}]
pub struct Operation {
    pub kind:   GateKind,     // typed gate enum
    pub qubits: Vec<usize>,   // qubit indices
    pub params: Vec<f64>,     // rotation angles
}
\end{lstlisting}

\subsection{Backends}

Four statevector backends are integrated:
\begin{enumerate}
  \item \textbf{Native} --- handwritten statevector simulator in Rust (reference implementation).
  \item \textbf{quantrs2} --- \texttt{quantrs2-core v0.2.0}, a Rust quantum simulation library.
  \item \textbf{roqoqo} --- HQS Quantum Simulations' \texttt{roqoqo v1.22.1} backend.
  \item \textbf{q1tsim} --- \texttt{q1tsim v0.5.0}, a gate-based Rust simulator.
\end{enumerate}

Additionally, \textsc{CleitonForge} provides a \textbf{DensityMatrixBackend}
(exact Kraus noise, up to 12 qubits) and a \textbf{NoisyStatevectorBackend}
(Monte Carlo wave-function trajectories).

% ─────────────────────────────────────────────────────────────────────────────
\section{The $R_z$ Sign Convention Divergence}

\subsection{Gate definition}

The $R_z(\lambda)$ gate is ubiquitous in quantum computing. IBM Quantum, Qiskit,
and the OpenQASM 3 standard library (\texttt{stdgates.inc}) define it as:

\begin{equation}
  R_z(\lambda)_{\mathrm{IBM}} =
  \begin{pmatrix} e^{-i\lambda/2} & 0 \\ 0 & e^{+i\lambda/2} \end{pmatrix}
  \label{eq:rz-ibm}
\end{equation}

The three backends native, \texttt{roqoqo}, and \texttt{q1tsim} all implement
\cref{eq:rz-ibm}. The \texttt{quantrs2-core} library uses the opposite sign:

\begin{equation}
  R_z(\lambda)_{\mathrm{reversed}} =
  \begin{pmatrix} e^{+i\lambda/2} & 0 \\ 0 & e^{-i\lambda/2} \end{pmatrix}
  \label{eq:rz-reversed}
\end{equation}

Both matrices are valid $R_z$ definitions---they differ by the global phase
convention of the computational basis. The choice is internally consistent
within each framework. However, when circuits are shared across frameworks
assuming a single convention, the mismatch becomes observable.

\subsection{Fidelity under the divergence}

Consider $\lambda = \pi/4$. Under \cref{eq:rz-ibm}: $R_z(\pi/4)\ket{0} =
e^{-i\pi/8}\ket{0}$. Under \cref{eq:rz-reversed}: $R_z(\pi/4)\ket{0} =
e^{+i\pi/8}\ket{0}$. The states are related by a global phase and remain
physically equivalent for single-qubit circuits.

In multi-qubit circuits, however, the phase becomes \emph{relative} when the
gate appears between entangling operations. For a QAOA cost layer
$e^{-i\gamma Z_0 Z_1}$ implemented as $\mathrm{CX} \cdot R_z(2\gamma) \cdot \mathrm{CX}$,
the effective two-qubit unitary differs between conventions. At $\gamma = -3\pi/4$:
the output states of the two conventions are orthogonal, yielding fidelity

\begin{equation}
  F = \abs{\braket{\psi_{\mathrm{IBM}}}{\psi_{\mathrm{reversed}}}}^2 = 0.
\end{equation}

% ─────────────────────────────────────────────────────────────────────────────
\section{Experimental Results}

\subsection{QAOA MaxCut benchmark}

We run the QAOA MaxCut circuit for the complete graph $K_2$ (2 qubits, 1 layer)
with $\gamma = -3\pi/4$, $\beta = -\pi/8$. Cross-backend statevector fidelity
$F = \abs{\braket{\psi_A}{\psi_B}}^2$ is measured for all pairs.

\begin{table}[h!]
\centering
\caption{Cross-backend fidelity on QAOA MaxCut ($K_2$, $\gamma=-3\pi/4$, $\beta=-\pi/8$).}
\label{tab:qaoa-fidelity}
\begin{tabular}{llrc}
\toprule
Backend A & Backend B & Fidelity & Status \\
\midrule
native   & roqoqo   & 1.00000000 & Agree \\
native   & q1tsim   & 1.00000000 & Agree \\
roqoqo   & q1tsim   & 1.00000000 & Agree \\
native   & quantrs2 & 0.00000000 & Diverge \\
roqoqo   & quantrs2 & 0.00000000 & Diverge \\
q1tsim   & quantrs2 & 0.00000000 & Diverge \\
\bottomrule
\end{tabular}
\end{table}

\subsection{Clifford benchmarks: no divergence}

Bell state, GHZ state, Bernstein--Vazirani, and Grover search all produce
cross-backend fidelity $F = 1.0$ for every pair of backends. These circuits
contain no $R_z$ gates at non-trivial angles, so the sign convention is
immaterial.

\subsection{Quantum Volume: also insensitive}

The IBM Quantum Volume benchmark~\cite{cross2019validating} evaluates simulators
using Haar-random $\mathrm{SU}(4)$ circuits. We compute the heavy-output
generation (HOG) fraction across 100 trials per circuit width.

\begin{table}[h!]
\centering
\caption{Quantum Volume HOG fractions (100 trials, seed fixed). All four backends
         agree to 4 decimal places. QV$= 2^n$ where HOG~$> 2/3$.}
\label{tab:qv}
\begin{tabular}{ccccccc}
\toprule
Width $n$ & QV & native & quantrs2 & roqoqo & q1tsim \\
\midrule
2 & ---    & 0.5355 & 0.5355 & 0.5355 & 0.5355 \\
3 & 8      & 0.7340 & 0.7340 & 0.7340 & 0.7340 \\
4 & 16     & 0.7991 & 0.7991 & 0.7991 & 0.7991 \\
5 & 32     & 0.8459 & 0.8459 & 0.8459 & 0.8459 \\
\bottomrule
\end{tabular}
\end{table}

\textbf{Why QV is insensitive.} Haar-random $\mathrm{SU}(4)$ circuits apply
$R_z$ at uniformly-distributed random angles. The sign flip $\lambda \to
-\lambda$ maps one random angle to another random angle drawn from the same
distribution. In expectation, the HOG fraction is unchanged. Only circuits
with \emph{specific} angles---such as QAOA's $\gamma$ and $\beta$---expose the
divergence.

\subsection{Performance}

All benchmarks run in release mode on a single core.

\begin{table}[h!]
\centering
\caption{Benchmark suite wall times (release build, single core).}
\label{tab:timing}
\begin{tabular}{lccr}
\toprule
Circuit & Qubits & Gates & Time (ms) \\
\midrule
Bell state                & 2 & 2  & 0.6 \\
GHZ state                 & 3 & 3  & 0.4 \\
QFT $\ket{100}$           & 3 & 8  & 0.4 \\
Bernstein--Vazirani       & 3 & 8  & 0.3 \\
Grover search             & 3 & 43 & 0.3 \\
QAOA MaxCut ($K_2$)       & 2 & 7  & 0.4 \\
\bottomrule
\end{tabular}
\end{table}

% ─────────────────────────────────────────────────────────────────────────────
\section{Convention Normalization Transpiler}

We implement a circuit transformation that converts between the two conventions
by negating angles on all $R_z$-family gates. The transformation is:

\begin{equation}
  \text{normalize}(R_z(\lambda)) = R_z(-\lambda)
\end{equation}

and analogously for $\text{Phase}(\lambda)$, $CR_z(\lambda)$, $CP(\lambda)$.
For the general unitary $U(\theta, \phi, \lambda)$, only $\phi$ and $\lambda$
are negated; $\theta$ is an $R_y$-style angle unaffected by the $R_z$
convention. For $CU(\theta, \phi, \lambda, \gamma)$, the global phase offset
$\gamma$ is also negated.

The transformation is implemented in \texttt{cforge\_parser::normalizer} and
exposed to Python as:

\begin{lstlisting}[style=python, caption={Convention normalization in Python.}]
import cforge

# Normalize from quantrs2 convention to IBM/Qiskit convention:
fixed = cforge.normalize(circuit, from_convention="reversed",
                                  to_convention="standard")

# Or directly from QASM source:
fixed = cforge.normalize_qasm(qasm_source)
\end{lstlisting}

After normalization, cross-backend fidelity between native and quantrs2
on the QAOA circuit is restored to $F = 1.00000000$.

% ─────────────────────────────────────────────────────────────────────────────
\section{Discussion}

\subsection{Implications for variational algorithms}

QAOA, VQE, ADAPT-VQE, and quantum machine learning circuits all rely on
$R_z$-family gates at angles determined by a classical optimizer. Users who
develop and optimize these circuits using \texttt{quantrs2} and then execute
them on IBM hardware, or compare with \texttt{roqoqo}-based simulators, will
silently obtain incorrect variational parameter landscapes. The optimization
will converge to the wrong parameters.

This class of error is particularly dangerous because:
\begin{enumerate}
  \item It is invisible to unit tests using Clifford gates.
  \item It is invisible to Quantum Volume and similar random-circuit benchmarks.
  \item The magnitude of the error is maximal ($F=0$) rather than small---there is no gradual degradation to warn the user.
\end{enumerate}

\subsection{Why the divergence exists}

The $R_z$ gate has two natural conventions depending on whether one considers
it a rotation of the qubit state or a rotation of the coordinate frame. The
IBM/Qiskit convention follows the standard physics rotation operator
$e^{-i\lambda Z/2}$. The reversed convention is $e^{+i\lambda Z/2}$, which is
the adjoint. Both are internally consistent and produce correct results for
circuits that are written and executed entirely within the same framework.

\subsection{Recommendation}

Framework authors should document their $R_z$ matrix definition explicitly in
the API reference, and ideally add a machine-readable metadata field in their
circuit representations. The OpenQASM standard~\cite{cross2022openqasm3}
specifies the IBM convention, and frameworks targeting hardware compatibility
should follow it.

Users exchanging circuits across frameworks should run a convention detection
test (a simple QAOA or single-qubit $R_z(\pi/4)$ fidelity check) before trusting
cross-framework results.

% ─────────────────────────────────────────────────────────────────────────────
\section{Related Work}

Cross-platform quantum circuit compatibility has been studied in the context of
hardware transpilation~\cite{qiskit2023}. The concept of neutral benchmarking
layers draws from classical compiler testing~\cite{yang2011finding}. The Quantum
Volume metric was introduced by Cross \emph{et al.}~\cite{cross2019validating}
as a hardware-independent performance metric.

To our knowledge, this is the first systematic multi-backend fidelity comparison
that explicitly identifies $R_z$ sign convention divergence as a source of
silent cross-backend disagreement in variational algorithm simulation.

% ─────────────────────────────────────────────────────────────────────────────
\section{Conclusion}

We have demonstrated that the $R_z$ gate sign convention is not universally
agreed upon among Rust quantum simulation libraries, and that this disagreement
produces catastrophic ($F=0$) fidelity loss on QAOA circuits while remaining
invisible to standard benchmarks. The \textsc{CleitonForge} neutral benchmarking
framework identified this discrepancy by routing identical circuits through four
independent backends using a canonical intermediate representation.

The normalization transpiler provided in \textsc{CleitonForge} allows users to
correct circuits for any target backend convention. Future work includes extending
the benchmark to additional frameworks, automating convention detection via backend
introspection, and extending the enterprise edition to provide per-device convention
calibration tables for hybrid classical-quantum workflows.

% ─────────────────────────────────────────────────────────────────────────────
\section*{Reproducibility}

All results are fully reproducible from the \textsc{CleitonForge} repository:

\begin{lstlisting}[style=rust]
git clone https://github.com/cleitonaugusto/cleitonforge
cd cleitonforge
cargo run --release --example benchmark_suite  -p cforge-cli
cargo run --release --example quantum_volume   -p cforge-cli
cargo run --release --example normalize_qaoa   -p cforge-cli
\end{lstlisting}

A Jupyter notebook reproducing all figures is at \texttt{notebooks/convention\_demo.ipynb}.

% ─────────────────────────────────────────────────────────────────────────────
\section*{Acknowledgments}

The author thanks the maintainers of quantrs2, roqoqo, and q1tsim for their
work on open-source quantum simulation infrastructure, which made this
comparison possible.

% ─────────────────────────────────────────────────────────────────────────────
\bibliographystyle{plain}
\begin{thebibliography}{9}

\bibitem{cross2019validating}
A.~W. Cross, L.~S. Bishop, S.~Sheldon, P.~D. Nation, and J.~M. Gambetta,
\emph{Validating quantum computers using randomized model circuits},
Physical Review A, 100:032328, 2019.

\bibitem{cross2022openqasm3}
A.~W. Cross, A.~Javadi-Abhari, T.~Alexander, N.~de~Beaudrap, L.~S. Bishop,
S.~Heidel, \emph{et al.},
\emph{OpenQASM 3: A broader and deeper quantum assembly language},
ACM Transactions on Quantum Computing, 3(3):1--50, 2022.

\bibitem{qiskit2023}
Qiskit contributors,
\emph{Qiskit: An open-source framework for quantum computing}, 2023.
\url{https://github.com/Qiskit/qiskit}

\bibitem{yang2011finding}
X.~Yang, Y.~Chen, E.~Eide, and J.~Regehr,
\emph{Finding and understanding bugs in C compilers},
In \emph{PLDI '11}, pages 283--294, 2011.

\end{thebibliography}

\end{document}
